\chapter{Normas y reglamentos}

El empleo de material radiactivo se rige por normas y reglamentos elaborados por la Comisión Nacional de Seguridad Nuclear y Salvaguardias (CNSNS). Para trabajar con material radiactivo se requiere:
\begin{itemize}
  \item una licencia que resuma los compromisos de la empresa con la CNSNS;
  \item autorización de la CNSNS para la adquisición de fuentes;
  \item notificación a la CNSNS para dar de baja fuentes.
\end{itemize}

A continuación se enumeran algunos requisitos de la CNSNS.

\section{Requisitos del personal ocupacionalmente expuesto}

El personal ocupacionalmente expuesto (POE) debe:
\begin{itemize}
  \item ser mayor de 18 años;
  \item tener escolaridad mínima secundaria;
  \item haber tomado un curso de seguridad radiológica;
  \item practicarse un examen médico de preempleo y después anualmente.
\end{itemize}

\section{Dosis, transporte y fuentes}

Respecto de las dosis se establece lo siguiente:
\begin{itemize}
  \item no recibir más de 5 rem al año;
  \item no exponer al público a dosis superiores a 500 mrem al año;
  \item informar al trabajador de su dosis acumulada anualmente.
\end{itemize}

Las unidades de transporte deben llevar como mínimo tres calcomanías y una caja portacontenedor.

Respecto de las fuentes se debe:
\begin{itemize}
  \item notificar las altas, bajas y localización geográfica;
  \item practicar pruebas de fuga semestralmente;
  \item mantener el contenedor con calcomanía y candado en buen estado;
  \item almacenarlas en lugares fijos, seguros y libres de material explosivo o inflamable;
  \item llevar la calcomanía de categoría de transporte I, II o III.
\end{itemize}

\section{Equipo de seguridad y obligaciones}

Todo el personal debe contar con dosímetros de bolsillo y de lectura directa. Debe haber una alarma y un detector por fuente; los detectores deben calibrarse semestralmente y los dosímetros verificarse también semestralmente.

Son obligaciones del POE seguir los procedimientos de la CNSNS establecidos en el manual y practicarse análisis clínicos anualmente.

Son obligaciones del permisionario mantener los archivos actualizados, renovar la licencia e informar a la CNSNS los movimientos de la fuente.

\section{Manual interno de seguridad radiológica}

El Manual Interno de Seguridad Radiológica contiene todas las políticas, criterios y procedimientos de seguridad radiológica. Sus disposiciones involucran al personal técnico, al personal administrativo, a los clientes y a las auditorías de la CNSNS.

\subsection{Límites administrativos y control de dosis del personal}

Para el POE A, personal técnico y auxiliar, se establece un límite mensual de $400\ \mathrm{mrem}$, menor que los 5 rem anuales establecidos por la CNSNS.

\subsection{Objetivos y control de dosis}

El objetivo es mantener controlada la dosis del personal. Se establecen los siguientes niveles mensuales:
\begin{itemize}
  \item nivel de investigación: $400\ \mathrm{mrem}$ acumulados al mes;
  \item nivel de intervención: $600\ \mathrm{mrem}$ acumulados al mes.
\end{itemize}

Para rebasar algún límite se debe llenar la solicitud SR-03, «Autorización de Dosis Adicional».

El control de las fuentes radiactivas corresponde únicamente al personal administrativo.

\subsection{Almacenamiento de fuentes radiactivas}

El almacén debe cumplir con las siguientes condiciones:
\begin{enumerate}[label=\alph*)]
  \item tener acceso controlado con llave;
  \item estar protegido de las condiciones ambientales;
  \item ser una instalación fija, no portátil;
  \item mantener un nivel de radiación en el exterior no mayor a $5\ \mathrm{mR/h}$;
  \item estar señalizado;
  \item no almacenar junto a él material explosivo o inflamable.
\end{enumerate}

\subsection{Transporte de fuentes radiactivas}

El vehículo contará con:
\begin{itemize}
  \item calcomanías de radiación;
  \item extintor;
  \item juego de herramientas;
  \item caja portacontenedor con candado;
  \item buenas condiciones de circulación.
\end{itemize}

Las fuentes no deben transportarse en compartimentos ocupados por pasajeros (autobús, ferrocarril, avión, correo, etcétera).

\subsection{Control del material y del equipo de seguridad radiológica}

El equipo se controla por residencias y mediante los formatos AL-01, «Resguardo de material y equipo», y AL-02, «Transferencia de material y equipo».

\subsection{Protección radiológica durante la toma de radiografías}

Antes y durante la toma de radiografías se debe:
\begin{itemize}
  \item verificar los dosímetros;
  \item verificar la alarma con el contenedor;
  \item verificar el detector, incluidas las baterías, la vigencia de la calibración y la respuesta;
  \item verificar el equipo de gammagrafía, asegurándose de que no tenga golpes ni estrangulamientos;
  \item acordonar 15 m si se utiliza gamma y 10 m si se utilizan rayos X;
  \item señalizar, manteniendo distancias no mayores de 10 m entre señales;
  \item vigilar el área mientras se esté exponiendo.
\end{itemize}

\subsection{Verificación y mantenimiento}

Al finalizar la exposición se debe:
\begin{itemize}
  \item verificar con el Geiger que la fuente regresó al contenedor;
  \item retirar las señales y los acordonamientos;
  \item no arrastrar las mangueras y desensamblar el equipo al trasladarse a otra área.
\end{itemize}

\subsection{Acordonamiento y señalización de áreas}

Los criterios de acordonamiento y señalización se describen en la sección anterior sobre protección radiológica durante la toma de radiografías.

\subsection{Uso y cuidado de los dosímetros}

Los dosímetros de lectura directa y de película:
\begin{itemize}
  \item se usarán siempre que se esté trabajando con radiación;
  \item se colocarán al frente, entre el cuello y la cintura;
  \item se vigilarán frecuentemente; el dosímetro de lectura directa se recargará cuando llegue a 100 mR o antes;
  \item no se debe humedecer el dosímetro de película;
  \item no se debe almacenar el dosímetro con la fuente.
\end{itemize}

\subsection{Operación y mantenimiento de los equipos}

La operación y el mantenimiento del equipo de gammagrafía deben efectuarse cada dos meses o antes si es necesario.

Para la operación y cuidado de los detectores de radiación y la alarma se debe:
\begin{itemize}
  \item evitar la humedad;
  \item evitar golpes y maltratos;
  \item protegerlos del polvo, aceite, grasa y otros contaminantes;
  \item almacenarlos cuando no se estén usando.
\end{itemize}

La calibración de los dosímetros de lectura directa y la calibración del detector Geiger Victoreen 492 corresponden únicamente al personal administrativo.

\section{Criterios a seguir en emergencias}

Se debe garantizar la seguridad física del operador y sus acompañantes, la integridad del contenedor y de la fuente, y la protección del público.

Siempre se debe llenar la forma SR-12, «Reporte de Emergencia Radiológica».

Se debe consultar el Manual para conocer las instrucciones específicas de cada accidente.
